\section*{Capítulo \ref{ch:g8:winds-global-circulation} --- Os ventos e a circulação geral}

\begin{solution}{exo:g8:winds-global-circulation:1}
(a) O alísio de nordeste, do nordeste. (b) Os \omterm{def:g8:winds-global-circulation:global-winds}{ventos de oeste}, do
noroeste. (c) Os \omterm{def:g8:winds-global-circulation:global-winds}{ventos polares de leste}, do leste (do nordeste). (d) O
alísio de sudeste, do sudeste.
\end{solution}

\begin{solution}{exo:g8:winds-global-circulation:2}
(a) O ar quente sobe e o ar frio desce. (b) Baixa. (c) Da alta pressão
para a baixa pressão. (d) Para a direita.
\end{solution}

\begin{solution}{exo:g8:winds-global-circulation:3}
(a) Falso: uma \omterm{def:g8:winds-global-circulation:breezes}{brisa marítima} sopra do \omterm{def:g2:water-on-earth:ocean}{mar} para a terra (de dia).
(b) Verdadeiro. (c) Falso: eles sopram do leste (do nordeste ou do
sudeste). (d) Verdadeiro.
\end{solution}

\begin{solution}{exo:g8:winds-global-circulation:4}
O desenho mostra um circuito: ao longo do chão, uma seta do \omterm{def:g2:water-on-earth:ocean}{mar} para a
terra; sobre a terra, uma seta subindo; lá no alto, uma seta da terra de
volta para cima do \omterm{def:g2:water-on-earth:ocean}{mar}; sobre o \omterm{def:g2:water-on-earth:ocean}{mar}, uma seta descendo. H fica marcado no
chão embaixo do ar que desce (sobre o \omterm{def:g2:water-on-earth:ocean}{mar}), L embaixo do ar que sobe
(sobre a terra).
\end{solution}

\begin{solution}{exo:g8:winds-global-circulation:5}
(a) 1.41 vez mais chão. (b) 50~\%. (c) 26~\%. (d) Perto dos \omterm{def:g3:continents-and-oceans:poles}{polos}, a luz
do Sol chega inclinada e se espalha por uma área grande, então cada metro
quadrado recebe pouca energia.
\end{solution}

\begin{solution}{exo:g8:winds-global-circulation:6}
(a) \omterm{def:g8:winds-global-circulation:global-winds}{Ventos polares de leste} (de \ang{70} a \ang{60} N), \omterm{def:g8:winds-global-circulation:global-winds}{ventos de oeste}
(de \ang{60} a \ang{30} N), \omterm{def:g8:winds-global-circulation:global-winds}{alísios} de nordeste (de \ang{30} N ao
\omterm{def:g3:continents-and-oceans:equator}{equador}), \omterm{def:g8:winds-global-circulation:global-winds}{alísios} de sudeste (do \omterm{def:g3:continents-and-oceans:equator}{equador} a \ang{30} S), \omterm{def:g8:winds-global-circulation:global-winds}{ventos de oeste}
(de \ang{30} a \ang{60} S), \omterm{def:g8:winds-global-circulation:global-winds}{ventos polares de leste} (de \ang{60} a
\ang{70} S). (b) Perto do \omterm{def:g3:continents-and-oceans:equator}{equador}, nas calmarias equatoriais. (c) Perto de
\ang{30} N e de \ang{30} S, embaixo do ar que desce.
\end{solution}

\begin{solution}{exo:g8:winds-global-circulation:7}
(a) A terra ($31 > 22$): o vento sopra do \omterm{def:g2:water-on-earth:ocean}{mar} para a terra, uma \omterm{def:g8:winds-global-circulation:breezes}{brisa
marítima}. (b) O \omterm{def:g2:water-on-earth:ocean}{mar} ($21 > 17$): o vento sopra da terra para o \omterm{def:g2:water-on-earth:ocean}{mar}, uma
\omterm{def:g8:winds-global-circulation:breezes}{brisa terrestre}. (c) A terra: $31 - 17 = 14$ graus; o \omterm{def:g2:water-on-earth:ocean}{mar}: $22 -
21 = 1$ grau. (d) Os raios do Sol só aquecem uma camada fina de chão, que
esquenta e esfria depressa; o \omterm{def:g2:water-on-earth:ocean}{mar} espalha o seu calor por uma camada
espessa de água, então a sua temperatura quase não muda entre o dia e a
noite.
\end{solution}

\begin{solution}{exo:g8:winds-global-circulation:8}
(a) Do norte. (b) Para a direita. (c) Do nordeste (é o alísio de
nordeste). (d) Sem a rotação, ele sopraria do sul; a rotação o desvia
para a esquerda; ele sopra do sudeste.
\end{solution}

\begin{solution}{exo:g8:winds-global-circulation:9}
Perto do \omterm{def:g3:continents-and-oceans:equator}{equador}, o chão e o \omterm{def:g2:water-on-earth:ocean}{mar} são fortemente \emph{aquecidos} pelo
Sol; o ar quente e úmido \emph{sobe}; enquanto sobe, ele \emph{esfria}, e
o seu vapor de água se transforma em \emph{\omterm{def:g4:clouds-rain-snow:cloud}{nuvens}}, que dão pancadas de
\omterm{def:g4:clouds-rain-snow:precipitation}{chuva} e trovoadas. Perto de \ang{25} N, o ar que subiu no \omterm{def:g3:continents-and-oceans:equator}{equador}
\emph{desce}; embaixo dele, a pressão no chão é \emph{alta}; o ar que
desce esquenta, as suas \omterm{def:g4:clouds-rain-snow:cloud}{nuvens} se desfazem e a região continua
\emph{árida}: um \omterm{def:g5:weather-and-climate:desert}{deserto}.
\end{solution}

\begin{solution}{exo:g8:winds-global-circulation:10}
(a) Os \omterm{def:g8:winds-global-circulation:global-winds}{ventos de oeste}. Não: eles sopram do oeste, contra um navio que
navega para o oeste. (b) Para chegar à faixa dos \omterm{def:g8:winds-global-circulation:global-winds}{alísios}, ao sul de cerca
de \ang{30} N. (c) O alísio de nordeste, entre o \omterm{def:g3:continents-and-oceans:equator}{equador} e cerca de
\ang{30} N. (d) Os \omterm{def:g8:winds-global-circulation:global-winds}{ventos de oeste}, ao norte de cerca de \ang{30} N (perto
de \ang{40} N, a \omterm{def:g6:latitude-longitude:latitude}{latitude} do porto de origem).
\end{solution}

\begin{solution}{exo:g8:winds-global-circulation:11}
(a) $1 \div 0.5 = 2$: a tabela dá 2. (b) $1 \div 0.866 \approx
1.1547$, ou seja, \textbf{1.15}. (c) A \ang{30}: $1000 \div 1.15 \approx
\qty{870}{W}$ em cada metro quadrado (exatamente $1000 \times 0.866 =
\qty{866}{W}$). A \ang{60}: $1000 \div 2 = \qty{500}{W}$.
\end{solution}

\begin{solution}{exo:g8:winds-global-circulation:12}
(a) Três circuitos: o ar sobe no \omterm{def:g3:continents-and-oceans:equator}{equador} e desce perto de \ang{30}; sobe
perto de \ang{60} e desce perto de \ang{30}; sobe perto de \ang{60} e desce
no \omterm{def:g3:continents-and-oceans:poles}{polo}. (b) L no \omterm{def:g3:continents-and-oceans:equator}{equador} e perto de \ang{60}; H perto de \ang{30} e no
\omterm{def:g3:continents-and-oceans:poles}{polo}. (c) De \ang{30} até o \omterm{def:g3:continents-and-oceans:equator}{equador}, o ar de superfície vai para o sul e
é desviado para a direita: alísio de nordeste. De \ang{30} a \ang{60},
ele vai para o norte e é desviado para a direita: ele vem do sudoeste. Do
\omterm{def:g3:continents-and-oceans:poles}{polo} até \ang{60}, ele vai para o sul, desviado para a direita: vento
polar de leste. (d) Na célula do meio, o ar de superfície corre em
direção ao \omterm{def:g3:continents-and-oceans:poles}{polo}, ou seja, para o norte; desviado para a direita, ele vai
para o nordeste: são os \omterm{def:g8:winds-global-circulation:global-winds}{ventos de oeste}.
\end{solution}

\begin{solution}{pb:g8:winds-global-circulation:1}
\textbf{1.} $59.6 - 23.5 = 36.1$ graus de \omterm{def:g6:latitude-longitude:longitude}{longitude}.

\textbf{2.} $36.1 \times 108 \approx 3899$: cerca de \qty{3900}{km}.

\textbf{3.} Sim: $3900 - 3896 = 4$~km, uma diferença minúscula.

\textbf{4.} Sim: os dois ficam perto de \ang{14} N, entre o \omterm{def:g3:continents-and-oceans:equator}{equador} e
cerca de \ang{30} N, a faixa do alísio de nordeste.

\textbf{5.} Do nordeste (as setas apontam para o sudoeste).

\textbf{6.} Para a frente, e um pouco para o lado: o vento vem de trás do
navio, um pouco da sua direita, e o empurra para o oeste.

\textbf{7.} Ele faz parte da \omterm{def:g8:winds-global-circulation:circulation}{circulação geral}: o ar continua descendo
perto de \ang{30} N e subindo perto do \omterm{def:g3:continents-and-oceans:equator}{equador}, então o ar de superfície
continua correndo em direção ao \omterm{def:g3:continents-and-oceans:equator}{equador}, e a rotação da Terra sempre o
desvia para a direita.

\textbf{8.} O alísio sopraria contra o navio. O capitão deveria navegar
para o norte, além de cerca de \ang{30} N, para encontrar os \omterm{def:g8:winds-global-circulation:global-winds}{ventos de
oeste}.

\textbf{9.} $7 \times 24 = \qty{168}{km}$ por dia.

\textbf{10.} $3 \times 168 = \qty{504}{km}$.

\textbf{11.} $3896 \div 168 \approx 23.2$ dias.

\textbf{12.} Cerca de \textbf{23 dias}: menos de um mês.
\end{solution}
